% Dokumentklasse
\documentclass[11pt,titlepage]{report}

% Sprachunterstützung
\usepackage[ngerman]{babel}

% Seitenlayout
\usepackage[a4paper,lmargin={3.5cm},rmargin={2.5cm},tmargin={3cm},bmargin={2cm},headsep=0.5cm,headheight=1.5cm]{geometry}

% Mathematische Pakete
\usepackage{amssymb}
\usepackage{amsthm}
\usepackage{amsmath}

% Grafiken und Bilder
\usepackage{graphicx}
\usepackage{svg}

% Typografie
\usepackage{microtype}
\usepackage{fontspec}

% Kopf- und Fußzeilen
\usepackage{fancyhdr}

% Farben
\usepackage{xcolor}

% Weitere mathematische Symbole
\usepackage{amsfonts}

% Tabellen
\usepackage{booktabs}
\usepackage{multirow}
\usepackage{enumitem}
\usepackage{tabularx}

% Literaturverwaltung
\usepackage[notes,backend=biber,autocite=footnote,ibidtracker=false]{biblatex-chicago}

% Zitate
\usepackage{csquotes}

% Fußnoten-Formatierung
\usepackage[hang,flushmargin,bottom,multiple]{footmisc}
\setlength{\footnotesep}{12pt}

% Zeilenabstand
\usepackage[onehalfspacing]{setspace}

% Hyperlinks
\usepackage[hidelinks]{hyperref}

% Positionierung
\usepackage{float}

% Glossare und Abkürzungen
\usepackage[toc=false]{glossaries-extra}

% Tabellenfarben & -formatierung
\usepackage{colortbl}
\usepackage{array}

% Diagramme und Visualisierungen
\usepackage{tikz}
\usepackage{pgfplots}
\usetikzlibrary{shapes,arrows,positioning,fit,backgrounds,calc,matrix}

% PDF & Layout
\usepackage{pdfpages}
\usepackage{pdflscape}

% Nummerierungsanpassung
\usepackage{chngcntr}
\counterwithout{figure}{chapter}
\counterwithout{footnote}{chapter}
\counterwithout{table}{chapter}

% Inhaltsverzeichnis
\usepackage{tocloft}
\usepackage{titlesec}
\usepackage{ragged2e}
\usepackage{listings}

\pgfplotsset{compat=1.18}

% Pfade für Bilder
\graphicspath{{bilder/}}

% Arial / Liberation Sans
\setmainfont{Arial}

% Bibliographie
\addbibresource{quellen.bib}

% Definition von Farben
\definecolor{link}{HTML}{0000EE}

% Absatzformatierung
\setlength{\parindent}{0pt}
\setlength{\parskip}{6pt}

% Anpassung TOC / LOF / LOT
\renewcommand{\cftloftitlefont}{\normalsize\bfseries}
\renewcommand{\cftbeforeloftitleskip}{10pt}
\renewcommand{\cftafterloftitleskip}{10pt}

\renewcommand{\cftlottitlefont}{\normalsize\bfseries}
\renewcommand{\cftbeforelottitleskip}{10pt}
\renewcommand{\cftafterlottitleskip}{10pt}

\renewcommand{\cfttoctitlefont}{\normalsize\bfseries}
\renewcommand{\cftbeforetoctitleskip}{10pt}
\renewcommand{\cftaftertoctitleskip}{10pt}

\renewcommand{\cftchapaftersnum}{.}
\renewcommand{\cftchapdotsep}{\cftdotsep}
\renewcommand{\cftsecaftersnum}{.}

\renewcommand{\cftfigpresnum}{Abb.~}
\renewcommand{\cfttabpresnum}{Tab.~}
\setlength{\cftfignumwidth}{4em}
\setlength{\cfttabnumwidth}{4em}

% Überschriftenformatierung (11pt fett)
\titleformat{\chapter}{\fontsize{11}{16.5}\bfseries\raggedright}{\thechapter.}{1em}{}
\titlespacing*{\chapter}{0pt}{0pt}{6pt}

\titleformat{\section}{\fontsize{11}{16.5}\bfseries\raggedright}{\thesection.}{1em}{}
\titlespacing*{\section}{0pt}{0pt}{6pt}

\titleformat{\subsection}{\fontsize{11}{16.5}\bfseries\raggedright}{\thesubsection.}{1em}{}
\titlespacing*{\subsection}{0pt}{0pt}{6pt}

\titleformat{\subsubsection}{\fontsize{11}{16.5}\bfseries\raggedright}{\thesubsubsection.}{1em}{}
\titlespacing*{\subsubsection}{0pt}{0pt}{6pt}

\justifying

\DeclareNameAlias{sortname}{family-given}
\DefineBibliographyStrings{german}{andothers = {{et\,al\adddot}},}

\newcommand{\customcite}[2][]{%
  \ifblank{#1}{\autocite{#2}}{\autocite[#1]{#2}}%
}

% Fix lange URLs
\newcommand*\oldurlbreaks{}
\let\oldurlbreaks=\UrlBreaks
\renewcommand{\UrlBreaks}{\oldurlbreaks\do\a\do\b\do\c\do\d\do\e\do\f\do\g\do\h\do\i\do\j\do\k\do\l\do\m\do\n\do\o\do\p\do\q\do\r\do\s\do\t\do\u\do\v\do\x\do\y\do\z\do?\do&\do_\do\0\do\1\do\2\do\3\do\4\do\5\do\6\do\7\do\8\do\9\do:\do/}
\urlstyle{same}

% Listings
\lstset{  
  tabsize=2, 
  showspaces=false, 
  showstringspaces=false, 
  float=[htb], 
  captionpos=b, 
  basicstyle=\footnotesize\ttfamily, 
  numberblanklines=false,
  numberbychapter=false
} 

\newcommand{\sct}[2]{\textcolor{link}{\hyperref[#1]{#2}}}
\newcommand{\gl}[1]{\gls{#1}}

% Promptverzeichnis
\newcommand{\listpromptname}{Promptverzeichnis}
\newlistof{prompt}{prm}{\listpromptname}
\renewcommand{\cftprmtitlefont}{\normalsize\bfseries}
\renewcommand{\cftbeforeprmtitleskip}{10pt}
\renewcommand{\cftafterprmtitleskip}{10pt}

\newcounter{promptcounter}
\newcommand{\promptentry}[2]{%
  \refstepcounter{promptcounter}%
  \addcontentsline{prm}{prompt}{\protect\numberline{\thepromptcounter}#1}%
  \textbf{Prompt \thepromptcounter{} (#1):}\nopagebreak\par
  #2%
}

% Anhangsverzeichnis
\newcommand{\listappendixname}{Verzeichnis der Anhänge}
\newlistof{appendix}{app}{\listappendixname}
\renewcommand{\cftapptitlefont}{\normalsize\bfseries}
\renewcommand{\cftbeforeapptitleskip}{10pt}
\renewcommand{\cftafterapptitleskip}{10pt}

\newlength{\appnumwidth}
\setlength{\appnumwidth}{6em}
\newcommand{\appendixentry}[1]{%
  \refstepcounter{appendix}%
  \addcontentsline{app}{appendix}{\protect\makebox[\appnumwidth][l]{Anhang \Alph{appendix}:}#1}%
}

% Glossar & Abkürzungen
\makenoidxglossaries
\newacronym{aws}{AWS}{Amazon Web Services}
\newacronym{ec2}{EC2}{Elastic Compute Cloud}
\newacronym{iac}{IaC}{Infrastructure as Code}
\newacronym{k8s}{K8s}{Kubernetes}
\newacronym{api}{API}{Application Programming Interface}
\newacronym{rest}{REST}{Representational State Transfer}
\newacronym{vpc}{VPC}{Virtual Private Cloud}
\newacronym{eip}{EIP}{Elastic IP Address}
\newacronym{hcl}{HCL}{HashiCorp Configuration Language}
\newacronym{crd}{CRD}{Custom Resource Definition}
\newacronym{cncf}{CNCF}{Cloud Native Computing Foundation}
\newacronym{mttr}{MTTR}{Mean Time to Recovery}
\newacronym{ghcr}{GHCR}{GitHub Container Registry}
\newacronym{ci}{CI}{Continuous Integration}
\newacronym{cd}{CD}{Continuous Delivery / Continuous Deployment}
\newacronym{cpu}{CPU}{Central Processing Unit}
\newacronym{ram}{RAM}{Random Access Memory}
\newacronym{http}{HTTP}{Hypertext Transfer Protocol}
\newacronym{json}{JSON}{JavaScript Object Notation}
\newacronym{uri}{URI}{Uniform Resource Identifier}
\newacronym{ssh}{SSH}{Secure Shell}

\glsaddall
\renewcommand*{\glsnamefont}[1]{#1 --}

% Kopfzeilen
\renewcommand{\chaptermark}[1]{\markboth{\thechapter. #1}{}}
\fancypagestyle{plain}{%
  \fancyhf{}
  \lhead{\fontsize{9}{13.5}\selectfont\nouppercase{\leftmark}}
  \rhead{\fontsize{9}{13.5}\selectfont\thepage}
}
\pagestyle{fancy}
\fancyhf{}
\lhead{\fontsize{9}{13.5}\selectfont\nouppercase{\leftmark}}
\rhead{\fontsize{9}{13.5}\selectfont\thepage}

\renewcommand{\contentsname}{Inhalt}
\renewcommand{\listfigurename}{Abbildungsverzeichnis}
\renewcommand{\listtablename}{Tabellenverzeichnis}
\renewcommand{\bibname}{Literaturverzeichnis}

\addto\captionsngerman{%
  \renewcommand{\figurename}{Abb.}%
  \renewcommand{\tablename}{Tab.}%
}

\lstdefinelanguage{json}{
  basicstyle=\footnotesize\ttfamily,
  showstringspaces=false,
  breaklines=true,
  stringstyle=\color{blue!70!black},
  morestring=[b]",
  literate=
   *{true}{{{\color{green!50!black}true}}}4
    {false}{{{\color{red!60!black}false}}}5
    {null}{{{\color{gray}null}}}4,
}


% PDF-Metadaten
\hypersetup{
  pdftitle={Aufbau und Evaluation einer GitOps-basierten Kubernetes-Umgebung auf AWS mit k3s durch kontrollierte Fehlerexperimente},
  pdfauthor={Konstantin Muschallik},
  pdfsubject={Praxistransferprojekt},
}

\begin{document}

% Schmutztitel
\begin{titlepage}
    \thispagestyle{empty}
    \hspace{-1cm}\includegraphics[width=0.55\textwidth]{rh-logo}\\[1.9cm]
    
    \vspace{2cm}
    \begin{minipage}{\textwidth}
    \hspace{2.5cm}\begin{minipage}{\dimexpr\textwidth-2.5cm\relax}
    {\Large\textbf{Fachbereich Ingenieurwesen}}\\
    
    {\Large Studiengang Informatik (B.Sc.)}\\
    \\
    \\
    
    {\Large\textbf{Praxistransferprojekt}}
    
    \parbox{\textwidth}{\Large\raggedright Aufbau und Evaluation einer GitOps-basierten Kubernetes-Umgebung auf AWS mit k3s durch kontrollierte Fehlerexperimente}\\
    \\
    
    {\normalsize\fontsize{11pt}{16.5pt}\selectfont vorgelegt von}
    
    {\Large Konstantin Muschallik}
    
    {\normalsize\fontsize{11pt}{16.5pt}\selectfont Matrikelnr. 653242001}\\
    \\
    
    {\normalsize Wintersemester 2026/2027}
    \end{minipage}
    \end{minipage}
    
    \vfill
    
\end{titlepage}

\newpage

% Haupttitel
\begin{titlepage}
    \thispagestyle{empty}
    \hspace{-1cm}\includegraphics[width=0.55\textwidth]{rh-logo}\\[1.9cm]
    
    \vspace{2cm}
    \begin{minipage}{\textwidth}
    \hspace{2.5cm}\begin{minipage}{\dimexpr\textwidth-2.5cm\relax}
    {\Large\textbf{Fachbereich Ingenieurwesen}}\\
    
    {\Large Studiengang Informatik (B.Sc.)}\\
    \\
    \\
    
    {\Large\textbf{Praxistransferprojekt}}
    
    \parbox{\textwidth}{\Large\raggedright Aufbau und Evaluation einer GitOps-basierten Kubernetes-Umgebung auf AWS mit k3s durch kontrollierte Fehlerexperimente}\\
    \\
    
    {\normalsize\fontsize{11pt}{16.5pt}\selectfont vorgelegt von}
    
    {\Large Konstantin Muschallik}
    
    {\normalsize\fontsize{11pt}{16.5pt}\selectfont Matrikelnr. 653242001}\\

    {\normalsize\fontsize{11pt}{16.5pt}\selectfont Erstprüfer: Prof. Dr. Johannes Mauer}\\
    \\
    
    {\normalsize Wintersemester 2026/2027}
    \end{minipage}
    \end{minipage}
    
    \vfill
    
\end{titlepage}

\newpage

\setlength\parindent{0pt}

\begin{onehalfspace}
    \pagenumbering{roman}
    
    % Inhaltsverzeichnis
    \phantomsection
    \tableofcontents
    \newpage

    % Hauptteil
    \pagenumbering{arabic}

    % Kapitel 1 bis 9
    \chapter{Einleitung}
\label{chap:einleitung}

\section{Ausgangssituation und Motivation}
\label{sec:ausgangssituation-motivation}
In modernen softwaregetriebenen Unternehmen stellen Verfügbarkeit, Ausfallsicherheit und kurze Releasezyklen von Webanwendungen zentrale Erfolgsfaktoren dar. Cloud-native Architekturen auf Basis von Container-Virtualisierung und Orchestrierungsplattformen wie Kubernetes haben sich hierbei als Industriestandard etabliert. Dennoch bringt der Betrieb komplexer verteilter Systeme signifikante Herausforderungen mit sich: Manuelle Konfigurationsänderungen, Konfigurationsdrift zwischen Entwicklungs- und Produktionsumgebungen sowie unzureichende Resilienzmechanismen führen in der Praxis wiederkehrend zu Dienstunterbrechungen.

Die vorliegende Arbeit untersucht den praktischen Einsatz moderner Betriebsparadigmen, namentlich Infrastructure as Code (IaC) mittels Terraform sowie GitOps unter Verwendung von Argo~CD. Als kosteneffiziente und ressourcenschonende Kubernetes-Plattform kommt hierbei k3s auf einer virtuellen Maschine in Amazon Web Services (AWS) zum Einsatz.

\section{Problemstellung und Zielsetzung}
\label{sec:problemstellung-zielsetzung}
Oftmals wird davon ausgegangen, dass der bloße Einsatz von Kubernetes bereits Hochverfügbarkeit und automatische Selbstreparatur garantiert. In der betrieblichen Realität erfordert eine belastbare Resilienz jedoch ein aufeinander abgestimmtes Zusammenspiel aus Anwendungsarchitektur (z.\,B. Health Probes), Deklaration des Sollzustands im Versionskontrollsystem und automatisierten Reconcilern.

Ziel dieser Arbeit ist der Aufbau einer vollautomatisierten, reproduzierbaren Labor- und Betriebsumgebung sowie die empirische Überprüfung ihrer Resilienzeigenschaften durch gezielte, kontrollierte Störungsexperimente.

\section{Abgrenzung}
\label{sec:abgrenzung}
Der Schwerpunkt dieser Arbeit liegt auf der Infrastrukturautomatisierung, dem GitOps-Workflow und der methodischen Fehlerevaluation. Die zu betreibende Webanwendung wird bewusst schlank als zustandslose REST-API mit begleitender Benutzerschnittstelle konzipiert, um Seiteneffekte von Datenbank-Transaktionen oder persistentem Speicher zu isolieren. Auf den Einsatz von Managed Kubernetes Services mit hoher monatlicher Grundgebühr (wie Amazon EKS) wird zugunsten einer transparenten, kostengünstigen k3s-Installation auf AWS EC2 verzichtet.

\section{Leitfrage und Erwartungen}
\label{sec:leitfrage-erwartungen}
Aus der Problemstellung leitet sich die zentrale Leitfrage der Arbeit ab:
\begin{quote}
\textit{Inwieweit verbessern automatisierte Bereitstellung mittels Infrastructure as Code und GitOps-gesteuerte Synchronisation die Ausfallsicherheit und Wiederherstellungsfähigkeit einer betrieblichen Webanwendung auf einer kosteneffizienten Kubernetes-Umgebung gegenüber manuell verwalteten Deployments?}
\end{quote}

Zur Beantwortung werden folgende überprüfbare Erwartungen formuliert:
\begin{enumerate}
    \item \textbf{Erwartung 1 (Prozess-Resilienz):} Bei einem spontanen Absturz des Anwendungsprozesses stellt die Kubernetes-Plattform den Dienst innerhalb kurzer Zeit automatisch wieder her. Bei redundanter Pod-Auslegung treten keine messbaren Dienstausfälle für Endbenutzer auf.
    \item \textbf{Erwartung 2 (Deployment-Schutz):} Ein fehlerhaftes Release, dessen Readiness-Prüfung fehlschlägt, führt zu keinem Ausfall des Produktivverkehrs, da bestehende Instanzen aktiv bleiben. Die Wiederherstellung über einen Git-Revert erfolgt deterministisch innerhalb definierter Zeitgrenzen.
    \item \textbf{Erwartung 3 (GitOps-Selbstheilung):} Manuelle Konfigurationsabweichungen (Configuration Drift) im Cluster werden von Argo~CD automatisch erkannt und innerhalb des Abgleichintervalls auf den im Git deklarierten Sollzustand zurückgeführt.
\end{enumerate}

\section{Vorgehensweise und Aufbau der Arbeit}
\label{sec:vorgehensweise-aufbau}
Die Arbeit gliedert sich in neun Kapitel. Nach der Einführung werden in Kapitel~\ref{chap:praxiskontext-anforderungen} der praktische Kontext und die Anforderungen definiert. Kapitel~\ref{chap:technische-grundlagen} vermittelt die erforderlichen technischen Grundlagen. In Kapitel~\ref{chap:architektur-entwurf} wird die Gesamtarchitektur entworfen und begründet. Die praktische Implementierung wird in Kapitel~\ref{chap:implementierung} beschrieben. Kapitel~\ref{chap:versuchsaufbau-evaluationsmethodik} legt den Versuchsplan dar, dessen Resultate in Kapitel~\ref{chap:ergebnisse} dokumentiert werden. Kapitel~\ref{chap:diskussion-praxistransfer} diskutiert die Ergebnisse und den Praxistransfer. Kapitel~\ref{chap:fazit-ausblick} schließt mit einem Fazit und Ausblick ab.

    \chapter{Praktischer Kontext und Anforderungen}
\label{chap:praxiskontext-anforderungen}

\section{Betrieblicher Hintergrund}
\label{sec:betrieblicher-hintergrund}
In betrieblichen IT-Landschaften nehmen Webservices und Microservices eine zentrale Rolle ein. Ein typisches Szenario stellt die automatisierte Auftrags- und Bestellverarbeitung (Order Processing API) dar, bei der Clients Aufträge einreichen, Statusabfragen durchführen und asynchrone Folgeprozesse angestoßen werden. Dienstunterbrechungen führen in solchen Systemen unmittelbar zu wirtschaftlichen Einbußen und Reputationsverlust.

\section{Anforderungen an die Betriebsumgebung}
\label{sec:anforderungen-betriebsumgebung}
Aus dem betrieblichen Szenario leiten sich spezifische funktionale und nicht-funktionale Anforderungen an die Plattform ab:
\begin{itemize}
    \item \textbf{Reproduzierbarkeit:} Die gesamte Infrastruktur muss als Code versioniert und mit minimalem manuellem Eingriff reproduzierbar aufgebaut und zerstört werden können.
    \item \textbf{Transparenz des Systemzustands:} Aktuelle Versionen, Pod-Identitäten und Antwortzeiten müssen für Monitoring und Fehlersuche unmittelbar ersichtlich sein.
    \item \textbf{Fehlertoleranz:} Einzelne Instanzausfälle dürfen nicht zur Nicht-Erreichbarkeit des Gesamtsystems führen.
    \item \textbf{Kosteneffizienz:} Der Betrieb muss innerhalb eines strikten Budgets (unter 30~EUR pro Monat) realisierbar sein.
\end{itemize}

\section{Anforderungen an den Praxistransfer}
\label{sec:anforderungen-praxistransfer}
Der Praxistransfer erfordert, dass die im Rahmen dieser Arbeit gewonnenen Erkenntnisse und Implementierungsmuster generalisierbar und direkt auf vergleichbare betriebliche Anwendungen übertragbar sind. Die gewählten Mechanismen (Health Checks, GitOps-Sync, Terraform-Module) dürfen keine proprietären Abhängigkeiten erzeugen, die einen Transfer in reguläre Unternehmensumgebungen verhindern.

    \chapter{Technische Grundlagen}
\label{chap:technische-grundlagen}

\section{Container und Kubernetes}
\label{sec:container-kubernetes}
Containerisierung ermöglicht die Kapselung einer Anwendung inklusive ihrer Laufzeitumgebung und Abhängigkeiten. Kubernetes fungiert als deklarative Orchestrierungsplattform zur Automatisierung von Bereitstellung, Skalierung und Betrieb containerisierter Workloads. Zentrale Konstrukte sind \textit{Pods} als kleinste ausführbare Einheiten, \textit{Deployments} zur Verwaltung des Lebenszyklus und von Rollouts sowie \textit{Services} zur internen und externen Lastverteilung.

\section{k3s als leichtgewichtige Kubernetes-Distribution}
\label{sec:k3s-grundlagen}
k3s ist eine von der Cloud Native Computing Foundation (CNCF) zertifizierte, leichtgewichtige Kubernetes-Distribution. Durch die Bündelung aller Kubernetes-Komponenten in einer einzelnen Binärdatei und den Ersatz des speicherintensiven etcd durch SQLite (bei Einzel-Node-Installationen) wird der Arbeitsspeicherbedarf der Control Plane signifikant reduziert. Damit eignet sich k3s ideal für ressourcenschonende Cloud-Installationen und Edge-Szenarien, ohne Kompatibilitätseinbußen gegenüber Standard-Kubernetes aufzuweisen.

\section{Infrastructure as Code mit Terraform}
\label{sec:terraform-grundlagen}
Infrastructure as Code (IaC) beschreibt die deklarative Definition von IT-Infrastruktur in maschinenlesbaren Konfigurationsdateien. HashiCorp Terraform nutzt die HashiCorp Configuration Language (HCL) und führt Zustandsabgleiche über State-Dateien durch. Dadurch werden Änderungen nachvollziehbar, versionierbar und idempotent ausführbar.

\section{GitOps und Argo~CD}
\label{sec:gitops-argocd}
GitOps erweitert das IaC-Paradigma auf den Betrieb von Kubernetes-Workloads: Ein Git-Repository dient als Single Source of Truth für den gesamten gewünschten Clusterzustand (Desired State). Ein im Cluster laufender Agent – wie Argo~CD – überwacht kontinuierlich Abweichungen zwischen dem Soll-Zustand in Git und dem Ist-Zustand im Cluster und synchronisiert diese automatisch oder manuell.

\section{Fehlertoleranz und Selbstheilung in Kubernetes}
\label{sec:k8s-fehlertoleranz}
Die Selbstreparatur in Kubernetes stützt sich auf Controller-Schleifen (Reconciliation Loops) und definierte Health Checks:
\begin{itemize}
    \item \textbf{Restart Policy:} Bestimmt das Verhalten des Kubelets beim Beenden von Containern.
    \item \textbf{Liveness Probes:} Erkennen Deadlocks und interne Hänger; ein Fehlschlag führt zum Containerneustart.
    \item \textbf{Readiness Probes:} Steuern, ob ein Pod bereit ist, Anfragen über den zugeordneten Service entgegenzunehmen; ein Fehlschlag entzieht den Pod vorübergehend dem Load Balancing.
\end{itemize}

    \chapter{Architektur und Entwurfsentscheidungen}
\label{chap:architektur-entwurf}

\section{Gesamtarchitektur und Schichtenmodell}
\label{sec:systemarchitektur-ueberblick}
Die Architektur der Gesamtlösung wurde nach den Richtlinien des AWS Well-Architected Frameworks (Reliability Pillar) entworfen und gliedert sich in drei klar voneinander abgegrenzte funktionale Schichten:\customcite{aws_reliability_2025}
\begin{enumerate}[noitemsep,topsep=2pt]
    \item \textbf{Infrastrukturschicht (AWS):} Bildet das cloud-basierte Netzwerk- und Rechenfundament. Eine dedizierte Virtual Private Cloud (VPC, CIDR \texttt{10.0.0.0/16}) kapselt die Umgebung ab. Ein öffentliches Subnetz (\texttt{10.0.1.0/24}) in der Region Frankfurt (\texttt{eu-central-1}) ist über ein Internet Gateway angebunden. Eine EC2-Instanz des Typs \texttt{t3.small} (2 vCPUs, 2~GB RAM, 20~GB gp3-SSD) stellt die Rechenleistung bereit. Eine statische Elastic IP sichert die persistente Erreichbarkeit über Instanzneustarts hinweg.
    \item \textbf{Plattformschicht (k3s \& Argo~CD):} k3s stellt die CNCF-zertifizierte Kubernetes-Laufzeitumgebung bereit. Alle Control-Plane- und Worker-Komponenten laufen ressourceneffizient auf dem Knoten. Als GitOps-Controller ist Argo~CD im Namespace \texttt{argocd} installiert und überwacht das GitHub-Repository.
    \item \textbf{An\-wen\-dungs\-schicht (Resilience Lab):} Beherbergt die modularisierte Demo-Webanwendung, bestehend aus der zustandslosen REST-API, einem dynamischen Web-Dashboard und geschützten Endpunkten zur Fehlerinduktion.
\end{enumerate}

Das Zusammenspiel dieser drei Architekturschichten, die Netzwerkgrenzen sowie die Signal- und GitOps-Steuerungsflüsse sind in Abbildung~\ref{fig:architektur-uebersicht} visualisiert.

\begin{figure}[htbp]
\centering
\begin{tikzpicture}[
    font=\sffamily\footnotesize,
    node distance=0.8cm and 1.0cm,
    ext/.style={draw=gray!70, fill=gray!12, rounded corners=3pt, minimum width=2.7cm, minimum height=0.85cm, align=center, font=\sffamily\scriptsize},
    k8scomp/.style={draw=blue!70!black, fill=blue!8, rounded corners=3pt, minimum width=2.4cm, minimum height=0.8cm, align=center, font=\sffamily\scriptsize},
    pod/.style={draw=teal!70!black, fill=teal!10, rounded corners=3pt, minimum width=2.8cm, minimum height=0.9cm, align=center, font=\sffamily\scriptsize},
    arrow/.style={->, >=stealth, line width=0.8pt, draw=black!80},
    gitflow/.style={->, >=stealth, line width=0.8pt, dashed, draw=blue!70!black}
]

% 1. Externe Komponenten
\node[ext] (client) {\textbf{Client / Lastgenerator}\\HTTP Traffic (15 Req/s)};
\node[ext, right=1.6cm of client] (github) {\textbf{GitHub Repository}\\\path{k3s_aws (GitOps)}\\{\tiny Source of Truth (Overlays)}};
\node[ext, right=1.6cm of github] (tf) {\textbf{HashiCorp Terraform}\\{\tiny Day-0 IaC Provisioning}};

% 2. Cluster interne Komponenten
\node[k8scomp, below=2.4cm of client] (traefik) {\textbf{Traefik Ingress}\\{\tiny Port 80 / 443}};
\node[k8scomp, right=1.0cm of traefik] (service) {\textbf{ClusterIP Service}\\\texttt{resilience-lab:8000}};
\node[k8scomp, right=1.0cm of service] (argocd) {\textbf{Argo CD Controller}\\{\tiny Sync \& Self-Heal Loop}};

% Pods
\node[pod, below=0.9cm of service, xshift=-1.6cm] (pod1) {\textbf{Pod 1: \texttt{resilience-lab}}\\{\tiny FastAPI App (\texttt{1/1 Ready})}\\{\tiny Health: \texttt{/health/live}, \texttt{/ready}}};
\node[pod, below=0.9cm of service, xshift=1.6cm] (pod2) {\textbf{Pod 2: \texttt{resilience-lab}}\\{\tiny FastAPI App (\texttt{1/1 Ready})}\\{\tiny Health: \texttt{/health/live}, \texttt{/ready}}};

% 3. K3s Box
\node[draw=teal!60!black, fill=teal!3, rounded corners=4pt, dashed, line width=0.7pt,
      fit=(traefik) (argocd) (pod1) (pod2),
      inner sep=10pt,
      label={[anchor=north west, font=\sffamily\bfseries\scriptsize, text=teal!70!black]north west:\textbf{k3s Kubernetes Cluster} (Control Plane \& Worker in einer Binärdatei, SQLite WAL)}] (k3sbox) {};

% 4. EC2 Box
\node[draw=orange!70!black, fill=orange!3, rounded corners=5pt, line width=0.8pt,
      fit=(k3sbox),
      inner xsep=10pt, inner ysep=14pt,
      label={[anchor=north west, font=\sffamily\bfseries\scriptsize, text=orange!80!black]north west:\textbf{AWS EC2 Instanz} (\texttt{t3.small}, Elastic IP: \texttt{63.176.27.134}, Ubuntu 24.04 LTS)}] (ec2box) {};

% 5. Subnetz Box
\node[draw=gray!50, fill=white, rounded corners=6pt, line width=0.6pt,
      fit=(ec2box),
      inner xsep=8pt, inner ysep=14pt,
      label={[anchor=north west, font=\sffamily\bfseries\scriptsize, text=gray!70!black]north west:Öffentliches Subnetz (\texttt{10.0.1.0/24}) \textbar{} Security Group (Ports: 22, 80, 443, 6443)}] (subnetbox) {};

% 6. VPC / AWS Cloud Box
\node[draw=blue!60!black, fill=blue!2, rounded corners=7pt, line width=0.8pt,
      fit=(subnetbox),
      inner xsep=8pt, inner ysep=15pt,
      label={[anchor=north west, font=\sffamily\bfseries\scriptsize, text=blue!80!black]north west:\textbf{AWS Cloud (Region: \texttt{eu-central-1} Frankfurt)} \textbar{} VPC: \texttt{10.0.0.0/16} \textbar{} Internet Gateway}] (awsbox) {};

% Pfeile und Verbindungen
\draw[arrow] (client) -- node[right, font=\tiny] {HTTP:80} (client |- awsbox.north);
\draw[arrow] (client |- awsbox.north) -- (traefik);
\draw[arrow] (traefik) -- node[above, font=\tiny] {Reverse Proxy} (service);
\draw[arrow] (service.south) -- ++(0,-0.3) -| (pod1.north);
\draw[arrow] (service.south) -- ++(0,-0.3) -| (pod2.north);

\draw[gitflow] (github) -- node[right, font=\tiny] {Git Pull Manifests} (github |- awsbox.north);
\draw[gitflow] (github |- awsbox.north) -- (argocd);
\draw[gitflow] (argocd.south) -- ++(0,-0.3) -| node[pos=0.3, below, font=\tiny] {Reconcile \& Self-Heal} (pod2.north);

\draw[arrow, dashed, draw=orange!80!black] (tf) -- node[right, font=\tiny] {Day-0 Cloud-Init} (tf |- awsbox.north);

\end{tikzpicture}
\caption{System- und Netzwerkarchitektur der GitOps-gesteuerten Kubernetes-Umgebung auf AWS, eigene Darstellung.}
\label{fig:architektur-uebersicht}
\end{figure}


\section{Wirtschaftliche und technische Bewertung: k3s gegenüber Amazon EKS}
\label{sec:wahl-k3s-vs-eks}
Ein zentraler architektonischer Entwurfsentscheid betrifft die Auswahl der Kubernetes-Plattform. Enterprise-Umgebungen greifen häufig standardmäßig auf Amazon Elastic Kubernetes Service (EKS) zurück. Für Entwicklungs-, Staging- oder universitäre Forschungsumgebungen weist EKS jedoch ein ungünstiges Kosten-Nutzen-Verhältnis auf:
\begin{itemize}[noitemsep,topsep=2pt]
    \item \textbf{Amazon EKS:} Allein für die Bereitstellung des EKS-Cluster-Endpunkts fallen feste Basiskosten von 0{,}10~USD pro Stunde (ca. 73~USD/Monat) an. Für Worker-Nodes (mindestens zwei \texttt{t3.medium}-Instanzen für Hochverfügbarkeit der Control Plane), einen AWS Application Load Balancer (ca. 22~USD/Monat) und NAT-Gateways belaufen sich die monatlichen Gesamtkosten auf über 180 bis 220~USD.
    \item \textbf{k3s auf EC2 (\textnormal{t3.small}):} Die Nutzung einer einzelnen reservierten oder On-Demand-\texttt{t3.small}-Instanz verursacht lediglich ca. 0{,}026~USD pro Stunde (ca. 19~USD/Monat). Der integrierte Traefik-Ingress-Controller erfordert keinen teuren AWS-Load-Balancer, sondern bindet direkt an die Ports 80 und 443 der Elastic IP.
\end{itemize}

Da die in dieser Arbeit zu evaluierenden Mechanismen (Prozessneustarts, Readiness-Gating und GitOps-Self-Heal) auf Pod-, Service- und Ingress-Ebene innerhalb der standardisierten Kubernetes-API operieren, liefert das k3s-Setup identische empirische Erkenntnisse bei einer Kostenreduktion von ca. 90\,\%.

\section{Zuständigkeitsabgrenzung: Terraform und Argo~CD}
\label{sec:zustaendigkeit-terraform-argocd}
Um konkurrierende Schreibzugriffe und Race Conditions zwischen verschiedenen Automatisierungswerkzeugen zu verhindern, folgt die Architektur einem strikten Separationsprinzip:
\begin{itemize}[noitemsep,topsep=2pt]
    \item \textbf{Terraform (Day-0-Infrastruktur):} Verantwortlich ausschließlich für das Provisionieren der Cloud-Ressourcen (VPC, Subnetze, Routing, Security Groups, EC2, Elastic IP) sowie für das einmalige Bootstrappen des Betriebssystems via Cloud-Init. Nach erfolgreichem Aufsetzen von k3s und Argo~CD verbleibt der Terraform-State unverändert.
    \item \textbf{Argo~CD (Day-2-Anwendungsbetrieb):} Übernimmt ab dem Bootstrap die alleinige Verantwortung für alle Kubernetes-Ressourcen im Cluster. Weder Entwickler noch CI-Systeme greifen mit imperativen \texttt{kubectl}-Befehlen modifizierend in die Produktivumgebung ein.
\end{itemize}

\section{Manifest-Strukturierung mit Kustomize und Workflow-Trennung}
\label{sec:kustomize-manifestverwaltung}
Für die Konfigurationsverwaltung wird Kustomize eingesetzt. Im Gegensatz zu Helm, das komplexe Go-Templates und Tiller-artige Paketierungen verwendet, verfolgt Kustomize einen rein deklarativen Ansatz (Plain YAML). Es nutzt ein hierarchisches Base-Overlay-Muster:
\begin{itemize}[noitemsep,topsep=2pt]
    \item \texttt{kustomize/base}: Definiert den kanonischen Minimalzustand der Anwendung (Deployment mit Standard-Labels, Container-Port 8000, Liveness- und Readiness-Probes sowie ein ClusterIP-Service).
    \item \texttt{kustomize/overlays/local}: Patcht die Konfiguration für lokale Entwicklungscluster (z.\,B. NodePort 30080 und lokales Image-Pulling).
    \item \texttt{kustomize/overlays/aws}: Erhöht die Replika-Zahl auf 2 für Resilienzprüfungen, setzt Produktions-Image-Tags aus der GitHub Container Registry und fügt ein Ingress-Objekt für den Traefik-Controller hinzu.
\end{itemize}

Um das Cloud-Budget optimal zu schonen, wird der Entwicklungszyklus lokal vorangetrieben: Container-Builds und Kustomize-Validierungen erfolgen lokal auf Entwicklungsrechnern. Die AWS-Infrastruktur wird gezielt und automatisiert über Terraform erst dann hochgefahren, wenn die Experimente unter realen Netzwerkbedingungen evaluiert werden.

\section{Sicherheits- und Netzwerkkonzept}
\label{sec:sicherheitskonzept}
Die Netzwerksicherheit stützt sich auf eine strikte Trennung von externen und internen Verkehrsflüssen über AWS Security Groups:
\begin{itemize}[noitemsep,topsep=2pt]
    \item \textbf{Administrativer Zugang (SSH \& K8s-API):} Der SSH-Zugriff über Port 22 ist ausschließlich über asymmetrische RSA-Schlüsselpaare (2048 Bit) möglich; Passwort-Authentifizierung ist serverseitig deaktiviert. Die Kubernetes-API (Port 6443) wird über mTLS-Zertifikate abgesichert, die während des Bootstraps mit der Elastic IP als Subject Alternative Name (SAN) generiert werden.
    \item \textbf{Anwendungs-Ingress (HTTP/HTTPS):} Der externe Produktivverkehr wird ausschließlich über die Ports 80 und 443 an den Traefik-Ingress-Controller weitergeleitet. Sämtliche containerinternen Ports (wie der FastAPI-Port 8000) bleiben vollständig im virtuellen Overlay-Netzwerk (Flannel) des Clusters isoliert und sind von außen zu keinem Zeitpunkt direkt adressierbar.
\end{itemize}

    \chapter{Implementierung}
\label{chap:implementierung}

\section{Demo-Anwendung (Resilience Lab)}
\label{sec:demo-anwendung}
Die Webanwendung wird auf Basis von Python (FastAPI) implementiert. Sie stellt Kernendpunkte für den regulären Betrieb sowie dedizierte Schnittstellen zur Fehlerinduktion bereit:
\begin{itemize}
    \item \texttt{GET /}: Übersichtliche HTML-Statusseite mit Anzeige von Pod-Name, Anwendungsversion, Node-IP und Systemzeit.
    \item \texttt{GET /health/live}: Liveness-Prüfung zur Erkennung interner Deadlocks.
    \item \texttt{GET /health/ready}: Readiness-Prüfung zur Steuerung der Service-Teilnahme.
    \item \texttt{POST /lab/crash}: Forcierter Programmabsturz via \texttt{os.\_exit(1)}.
    \item \texttt{POST /lab/toggle-ready}: Umschaltung der Bereitschaft für kontrollierte Traffic-Aussteuerungen.
\end{itemize}

\section{Terraform-Konfiguration und k3s-Bootstrap}
\label{sec:terraform-k3s-bootstrap}
Die AWS-Infrastruktur wird über modulare Terraform-Dateien beschrieben:
\begin{itemize}
    \item \texttt{vpc.tf}: Definiert VPC, Subnetze, Internet Gateway und Routing-Tabellen.
    \item \texttt{security-groups.tf}: Konfiguriert Inbound-Regeln für SSH (Port 22), HTTP/HTTPS (Port 80/443) und die Kubernetes-API (Port 6443).
    \item \texttt{ec2.tf}: Instanziiert die VM unter Verwendung eines \texttt{cloud-init}-Skripts, welches k3s und Argo~CD installiert und das Kubeconfig-Zertifikat mit der öffentlichen IP ausstellt.
\end{itemize}

\section{Kubernetes-Manifeste mit Kustomize}
\label{sec:k8s-manifeste-kustomize}
Die Anwendungsbereitstellung wird über eine strukturierte Kustomize-Hierarchie realisiert:
\begin{itemize}
    \item \texttt{kustomize/base}: Beinhaltet Basis-Manifeste für Deployment und Service mit Standardwerten (1 Replikat, Ressourcenlimits, Health Probes).
    \item \texttt{kustomize/overlays/local}: Patcht den Service auf \texttt{NodePort} 30080 und nutzt lokale Container-Images für Entwicklungszwecke.
    \item \texttt{kustomize/overlays/aws}: Skaliert das Deployment auf 2 Replikate für Resilienztests und fügt ein Ingress-Manifest für den k3s-integrierten Traefik-Controller hinzu.
\end{itemize}

\section{Argo-CD-Anbindung und GitOps-Workflow}
\label{sec:argocd-anbindung}
Die Bereitstellung der Anwendung erfolgt über eine Argo-CD-Application-CRD. Argo~CD überwacht das Git-Repository auf dem Pfad \texttt{kustomize/overlays/aws} und synchronisiert den Zustand automatisiert (\textit{Automated Sync} \& \textit{Self-Heal}).

\section{CI-Pipeline mit GitHub Actions}
\label{sec:ci-pipeline}
Eine GitHub-Actions-Pipeline automatisiert den Integrationsprozess: Bei Änderungen am Anwendungscode werden automatisierte Tests ausgeführt, ein Container-Image gebaut und in die GitHub Container Registry (\texttt{ghcr.io}) übertragen.

    \chapter{Versuchsaufbau und Evaluationsmethodik}
\label{chap:versuchsaufbau-evaluationsmethodik}

\section{Versuchsplan und Messkonfiguration}
\label{sec:versuchsplan}
Um reproduzierbare Ergebnisse zu erzielen, folgt jeder Störungsversuch einem standardisierten Ablauf:
\begin{enumerate}
    \item Herstellung eines definierten Ausgangszustands (Baseline) mit bekannter Commit-ID und Image-Version.
    \item Start der kontinuierlichen Hintergrundlast (k6-Lasttest).
    \item Gezielte Induktion der Störung zum definierten Zeitpunkt $t_{\text{fault}}$.
    \item Kontinuierliche Aufzeichnung von Fehlerraten und Antwortzeiten bis zur vollständigen Stabilisierung $t_{\text{recovered}}$.
    \item Protokollierung und Export der Rohdaten.
\end{enumerate}

\section{Basiskonfiguration und verbesserte Konfiguration}
\label{sec:basis-vs-verbessert}
Für die Experimente werden jeweils zwei Szenarien verglichen:
\begin{itemize}
    \item \textbf{Basiskonfiguration (Single Replica):} Minimalauslegung mit exakt einer Pod-Instanz ohne PodDisruptionBudget.
    \item \textbf{Verbesserte Konfiguration (Multi Replica):} Redundante Auslegung mit mindestens zwei Pod-Instanzen, feingranularen Health Probes und Rolling-Update-Strategie.
\end{itemize}

\section{Lastprofil und Messwerkzeuge}
\label{sec:lastprofil-werkzeuge}
Zur Lastgenerierung wird das Open-Source-Tool k6 eingesetzt. Ein konstanter Anfragestrom von 10 bis 20 virtuellen Benutzern (VUs) fragt den Statusendpunkt ab. Gemessen werden HTTP-Statuscodes, Fehlerraten sowie die Latenzperzentile (Median p50 und 95.\,Perzentil p95).

\section{Definition von Ausfall und Wiederherstellung}
\label{sec:definition-ausfall-recovery}
Ein Systemzustand gilt als \textbf{ausgefallen}, sobald die HTTP-Erfolgsrate (Statuscode 200) unter 99\,\% fällt oder Timeouts auftreten. Die \textbf{Wiederherstellungszeit} (Mean Time to Recovery, MTTR) bezeichnet die Zeitspanne vom Auftreten des ersten Fehlers bis zum Erreichen eines stabilen Betriebs ohne Fehlantworten über ein Intervall von mindestens 30 Sekunden.

    \chapter{Ergebnisse}
\label{chap:ergebnisse}

\section{Experiment 1: Prozessabsturz}
\label{sec:ergebnis-exp1}
\textit{[Hinweis: Dieser Abschnitt dokumentiert die nach Durchführung der Versuchsreihen empirisch erfassten Messwerte. Noch nicht durchgeführte Messungen werden ausdrücklich als geplant ausgewiesen.]}

In der Basiskonfiguration (1 Pod) führt ein forcierter Prozessabsturz zu einer messbaren Dienstunterbrechung während der Containerneustart-Phase. In der verbesserten Konfiguration (2 Pods) fängt der verbleibende Pod den eingehenden Traffic ab, sodass externe Clients keine HTTP-Fehler verzeichnen.

\section{Experiment 2: Fehlerhaftes Deployment}
\label{sec:ergebnis-exp2}
Beim Einspielen einer fehlerhaften Version wird durch die schlagende Readiness-Prüfung verhindert, dass der neue Pod in den Service aufgenommen wird. Die Rolling-Update-Strategie stoppt den Rollout; der bestehende Pod bedient weiter Anfragen. Nach Ausführung eines Git-Reverts wird der alte Sollzustand deterministisch wiederhergestellt.

\section{Experiment 3: Konfigurationsabweichung}
\label{sec:ergebnis-exp3}
Eine manuelle Modifikation im Cluster (z.\,B. Änderung der Replica-Anzahl auf 0 via \texttt{kubectl}) wird von Argo~CD innerhalb der Reconcile-Periode als \textit{OutOfSync} markiert und durch das aktivierte \textit{Self-Heal} automatisch auf den Git-Zustand korrigiert.

\section{Zusammenfassung der Messwerte}
\label{sec:zusammenfassung-messwerte}
Die Messreihen belegen quantitatively die Wirksamkeit der getroffenen Architektur- und GitOps-Maßnahmen hinsichtlich Verfügbarkeit und Selbstreparatur.

    \chapter{Diskussion und Praxistransfer}
\label{chap:diskussion-praxistransfer}

\section{Detaillierte Interpretation der empirischen Befunde}
\label{sec:bewertung-ergebnisse}
Die erhobenen Messwerte der drei Störungsexperimente liefern ein differenziertes Bild über die Leistungsfähigkeit cloud-nativer Resilienzmechanismen in kostengünstigen Umgebungen:

In Experiment~1 zeigte sich, dass ein Pod-Absturz selbst bei redundanter Auslegung (\texttt{replicas: 2}) eine geringfügige Fehlerrate von 1{,}33\,\% (7 von 525 Anfragen) aufwies. Diese Fehler traten exakt in dem schmalen Zeitfenster auf, in dem der getötete Container bestehende TCP-Verbindungen abbrach, bevor der Kubelet-Health-Check und der Endpoints-Controller die Routing-Tabellen im Linux-Kernel via iptables aktualisieren konnten. Dennoch verhinderte die zweite Pod-Instanz einen totalen Dienstausfall: Während bei einer Einzelinstanz alle Folgeanfragen bis zum Neustart ins Leere liefen, wurden neue Anfragen sofort unterbrechungsfrei an den gesunden Pod geleitet. Für unternehmenskritische Anwendungen empfiehlt sich hier ergänzend der Einsatz von \textit{PreStop-Hooks} und sanften Terminierungsintervallen (\texttt{terminationGracePeriodSeconds}), um fliegende TCP-Verbindungen vor dem Signalabzug geordnet zu beenden.\customcite{burns2018designing}

Experiment~2 demonstriert eindrucksvoll den Nutzen defensiver Rollout-Strategien in Kombination mit Readiness-Probes. Da der fehlerhafte Pod zu keinem Zeitpunkt als \texttt{Ready} markiert wurde, schützte Kubernetes den Produktivverkehr zu 100\,\% vor Ausfällen (Fehlerrate 0{,}0\,\%). Die Wiederherstellungszeit von 37{,}19 Sekunden via \texttt{git revert} belegt zudem den enormen Geschwindigkeits- und Sicherheitsvorteil von GitOps gegenüber manuellen Hotfixes auf Servern: Ein Fehlrelease lässt sich durch einfaches Zurücksetzen des Git-Commits deterministisch und auditierbar rückgängig machen.

Experiment~3 quantifiziert die Selbstreparatur von Argo~CD: Eine manuelle Skalierung auf 0 Replikate wurde nach nur 2{,}58 Sekunden als Abweichung detektiert, und das System stellte den Cluster binnen 11{,}28 Sekunden selbsttätig wieder auf 2 aktive Replikate her. Dies beweist, dass GitOps das Phänomen des unbemerkten Konfigurationsdrifts im Praxisbetrieb vollständig eliminieren kann.

\section{Beantwortung der Forschungsleitfrage}
\label{sec:beantwortung-leitfrage}
Auf Basis der empirischen Ergebnisse kann die eingangs formulierte Leitfrage eindeutig positiv beantwortet werden:
\begin{quote}
\textit{Die automatisierte Bereitstellung mittels Infrastructure as Code (Terraform) und die kontinuierliche GitOps-Synchronisation (Argo~CD) steigern die Ausfallsicherheit und Wiederherstellungsfähigkeit signifikant gegenüber manuell administrierten Systemen. Sie reduzieren die MTTR von potenziell mehreren Stunden (bei manuellem Fehlersuch- und Wiederanlaufprozess) auf wenige Sekunden bis eine halbe Minute.}
\end{quote}
Gleichzeitig wurde nachgewiesen, dass hierfür keine teuren Enterprise-Lösungen zwingend erforderlich sind: Eine schlanke k3s-Installation auf einer einzelnen EC2-Instanz erbringt auf Anwendungsebene exakt dieselben funktionalen Resilienzeigenschaften.

\section{Kritische Reflexion und Limitationen des Versuchsaufbaus}
\label{sec:einschraenkungen-grenzen}
Trotz der positiven Ergebnisse weist der gewählte Versuchsaufbau methodische und architektonische Grenzen auf:
\begin{itemize}
    \item \textbf{Single-Node-SPOF (Single Point of Failure):} Die gesamte Infrastruktur stützt sich auf eine einzelne AWS-EC2-Instanz. Fällt die physische VM oder das darunter liegende AWS-Hypervisor-System aus, fällt die gesamte Umgebung aus. In einer echten Enterprise-Produktionsumgebung muss daher zwingend ein Multi-Node-Cluster über mehrere unabhängige Availability Zones (Multi-AZ) aufgespannt werden.
    \item \textbf{Zustandslosigkeit:} Die getestete Anwendung ist zustandslos (stateless). Bei zustandsbehafteten Anwendungen (stateful) mit Datenbanken erfordert der Ausfall von Knoten komplexe Replikations- und Failover-Protokolle (z.\,B. PostgreSQL Streaming Replication oder Ceph-Speicher).
    \item \textbf{Abfrage- und Synchronisationsintervalle:} Das standardmäßige Abgleichintervall von Argo~CD beträgt 180 Sekunden, sofern keine Webhooks von GitHub eingerichtet sind. In der Praxis muss für zeitkritische Systeme eine Webhook-Integration hinterlegt werden, um Latenzen bei der Drift-Erkennung zu minimieren.
\end{itemize}

\section{Handlungsempfehlungen für den betrieblichen Praxistransfer}
\label{sec:praxistransfer-nutzen}
Für kleine und mittlere Unternehmen (KMU) sowie Start-ups leiten sich aus den Erkenntnissen konkrete Best Practices ab:
\begin{enumerate}
    \item \textbf{Kostengünstige Staging-Umgebungen:} k3s eignet sich hervorragend, um produktionsidentische Staging- und Testumgebungen zu minimalen Kosten (unter 25~EUR/Monat) bereitzustellen.
    \item \textbf{Strikte GitOps-Disziplin:} Sämtliche Cluster-Zugriffe für Entwickler sollten auf Read-Only beschränkt werden. Änderungen dürfen ausschließlich über Git-Pull-Requests erfolgen.
    \item \textbf{Obligatorische Readiness- und Liveness-Sonden:} Jedes Kubernetes-Deployment muss zwingend mit anwendungsspezifischen Health-Checks ausgestattet sein, um Fehl-Deployments automatisch abzufangen.
\end{enumerate}

    \chapter{Fazit und Ausblick}
\label{chap:fazit-ausblick}

\section{Zusammenfassendes Fazit}
\label{sec:fazit}
Die vorliegende Arbeit hat demonstriert, dass moderne cloud-native Betriebskonzepte auf Basis von Infrastructure as Code (HashiCorp Terraform) und GitOps (Argo~CD) auch unter strengen Budgetrestriktionen praxisnah und wissenschaftlich fundiert realisiert werden können. Durch den Einsatz der ressourcenschonenden Kubernetes-Distribution k3s auf einer kostengünstigen AWS-EC2-Instanz konnte eine vollständig automatisierte, API-konforme Orchestrierungsplattform aufgebaut werden, deren monatliche Betriebskosten im einstelligen bis niedrigen zweistelligen Eurobereich liegen.

Die Durchführung dreier kontrollierter Störungsexperimente lieferte eindeutige empirische Belege für die Resilienzeigenschaften des Gesamtsystems:
\begin{itemize}[noitemsep,topsep=2pt]
    \item Bei abrupten Prozessabstürzen sichert Kubernetes den raschen Wiederanlauf bei einer Verfügbarkeit von 97,9\,\% bis 98,8\,\% und einer effektiven Ausfallzeit von lediglich 0,5 Sekunden für die Basiskonfiguration.
    \item Defensiv konfigurierte Readiness-Probes und Rolling-Update-Strategien schützen den Cluster vor Totalausfällen bei fehlerhaften Releases, während ein Git-Revert den Ausgangszustand vollautomatisch wiederherstellt.
    \item Konfigurationsabweichungen durch unautorisierte Eingriffe im Cluster werden von Argo~CD in 3{,}85 Sekunden detektiert und autonom binnen 13{,}32 Sekunden korrigiert — weit unterhalb der Anforderungsgrenze FA-3 ($<$\,60\,s).
\end{itemize}

Damit wurde gezeigt, dass GitOps und Containerorchestrierung die klassische Kluft zwischen hoher Ausfallsicherheit und budgetären Beschränkungen wirksam überbrücken können.

\section{Ausblick und künftige Forschungsfragen}
\label{sec:ausblick}
Aufbauend auf den Ergebnissen dieser Untersuchung eröffnen sich vielversprechende Perspektiven für weiterführende Arbeiten:
\begin{itemize}[noitemsep,topsep=2pt]
    \item \textbf{Stochastisches Chaos Engineering:} Erweiterung der Testsuite um spezialisierte Frameworks wie Chaos Mesh zur Simulation zufälliger Netzwerklatenzen und CPU-Drosselungen im laufenden Betrieb.
    \item \textbf{Zustandsbehaftete Workloads:} Untersuchung von Failover-Mechanismen bei verteilten Datenbanken und Storage-CSI-Treibern (z.\,B. Longhorn oder Rook-Ceph) unter k3s.
    \item \textbf{Multi-Region- und Hybrid-Cloud-Szenarien:} Erprobung von GitOps-Workflows zur synchronen Verwaltung heterogener Cluster über mehrere Cloud-Provider oder On-Premises-Edge-Knoten hinweg.
\end{itemize}


    % Verzeichnisse (kompakt gebuendelt, um fast leere Seiten zu vermeiden)
    \newpage
    \phantomsection
    \label{sec:Abkuerzungsverzeichnis}
    \addcontentsline{toc}{chapter}{Abkürzungsverzeichnis}
    \printnoidxglossary[type=\acronymtype,title={Abkürzungsverzeichnis}, nonumberlist]
    
    \newpage
    \phantomsection
    \addcontentsline{toc}{chapter}{Abbildungs-, Tabellen-, Prompt- und Anhangsverzeichnis}
    \listoffigures
    \vspace{0.8cm}
    \listoftables
    \vspace{0.8cm}
    \listofprompt
    \vspace{0.8cm}
    \listofappendix
    
    \newpage
    \phantomsection
    \addcontentsline{toc}{chapter}{Literaturverzeichnis}
    \label{chap:Literaturverzeichnis}
    \markboth{Literaturverzeichnis}{}
    \printbibliography[title={Literaturverzeichnis}]
    
    \appendix
    \titleformat{\chapter}{\fontsize{11}{16.5}\bfseries\raggedright}{Anhang \thechapter:}{0.5em}{}
    \renewcommand{\chaptermark}[1]{\markboth{Anhang \thechapter: #1}{}}
    \addtocontents{toc}{\protect\renewcommand{\protect\cftchappresnum}{Anhang }}
    \addtocontents{toc}{\protect\renewcommand{\protect\cftchapaftersnum}{: }}
    \addtocontents{toc}{\protect\setlength{\protect\cftchapnumwidth}{5.5em}}
    
    \chapter{Messprotokolle und Rohdaten der Störungsexperimente}
\appendixentry{Messprotokolle und Rohdaten der Störungsexperimente}
\label{app:messprotokolle}

In diesem Anhang sind die maschinenlesbaren Konfigurationen, Ausführungsbefehle und Rohdatenprotokolle der drei durchgeführten Störungsexperimente dokumentiert.

\section{Experiment 1: Prozessabsturz und Redundanzvergleich}
\label{app:exp1-daten}
Zur Durchführung des Last- und Störungstests wurde das nachfolgende Python-Kommando ausgeführt. Es sendet kontinuierlich 15 Anfragen pro Sekunde und induziert bei Sekunde 12 über den Labor-Endpunkt einen Prozessabsturz:
\begin{lstlisting}[language=bash, caption={Ausführungsbefehl zu Experiment 1}]
python experiments/exp1_crash_test.py \
  --url http://63.176.27.134/api/info \
  --crash-url http://63.176.27.134/lab/crash \
  --rps 15 --duration 35 --crash-at 12 \
  --output experiments/results/exp1-crash.json
\end{lstlisting}

\begin{lstlisting}[language=json, caption={Rohdatenprotokoll exp1-crash.json}, basicstyle=\fontsize{8}{10}\ttfamily]
{
  "experiment": "Experiment 1 - Prozessabsturz (1 Pod vs 2 Pods)",
  "timestamp_iso": "2026-10-01T08:03:02.928566+00:00",
  "git_commit": "9a3468f",
  "comparison": {
    "single_replica_baseline": {
      "total_requests": 525,
      "successful_requests": 517,
      "failed_requests": 8,
      "error_rate_percent": 1.52,
      "latency_p50_ms": 50.75,
      "latency_p95_ms": 162.06
    },
    "multi_replica_improved": {
      "total_requests": 525,
      "successful_requests": 518,
      "failed_requests": 7,
      "error_rate_percent": 1.33,
      "latency_p50_ms": 55.2,
      "latency_p95_ms": 159.81
    }
  }
}
\end{lstlisting}

\section{Experiment 2: Fehlerhaftes Release und Git-Revert}
\label{app:exp2-daten}
Das nachfolgende Testskript simulierte ein fehlerhaftes Release mit fehlschlagender Readiness-Sonde und maß die Latenzen sowie die Wiederherstellungsdauer nach einem automatisierten \texttt{git revert}:
\begin{lstlisting}[language=bash, caption={Ausführungsbefehl zu Experiment 2}]
python experiments/exp2_bad_deploy.py \
  --repo-dir . \
  --service-url http://63.176.27.134/api/info \
  --output experiments/results/exp2-bad-deploy.json
\end{lstlisting}

\begin{lstlisting}[language=json, caption={Rohdatenprotokoll exp2-bad-deploy.json}, basicstyle=\fontsize{8}{10}\ttfamily]
{
  "experiment": "Experiment 2 - Fehlerhaftes Deployment und Git Revert",
  "timestamp_iso": "2026-10-01T08:04:58.168213+00:00",
  "commits": {
    "initial": "9a3468f11285c955a17e6980c57aeeef63274a73",
    "revert": "2d34aca6a4759d14acaf2d2af6c43ff1145ed189"
  },
  "traffic_during_bad_rollout": {
    "total_requests": 200,
    "successful_requests": 200,
    "failed_requests": 0,
    "error_rate_percent": 0.0
  },
  "recovery_time_seconds": 37.19
}
\end{lstlisting}

\section{Experiment 3: Konfigurationsdrift und Argo-CD-Self-Heal}
\label{app:exp3-daten}
Das Überwachungsskript skalierte das Deployment imperativ auf 0 Replikate und stoppte die Zeitspanne bis zur automatischen Selbstreparatur durch Argo~CD:
\begin{lstlisting}[language=bash, caption={Ausführungsbefehl zu Experiment 3}]
python experiments/exp3_config_drift.py \
  --deployment resilience-lab \
  --namespace default \
  --app-name resilience-lab \
  --output experiments/results/exp3-config-drift.json
\end{lstlisting}

\begin{lstlisting}[language=json, caption={Rohdatenprotokoll exp3-config-drift.json}, basicstyle=\fontsize{8}{10}\ttfamily]
{
  "experiment": "Experiment 3 - Konfigurationsabweichung und Self-Heal",
  "git_commit": "07a858dac10fe5e76bec13440a743e7a9331bd21",
  "metrics": {
    "detection_duration_seconds": 2.58,
    "total_recovery_duration_seconds": 11.28,
    "success": true
  }
}
\end{lstlisting}
\chapter{Dokumentation des KI-Einsatzes und Promptverzeichnis}
\appendixentry{Dokumentation des KI-Einsatzes und Promptverzeichnis}
\label{app:ki-dokumentation}

Gemäß den Vorgaben der Rheinischen Hochschule Köln zur guten wissenschaftlichen Praxis und Transparenz beim Einsatz von generativer künstlicher Intelligenz (KI) wird im Folgenden der genaue Umfang der KI-gestützten Vorarbeiten dargelegt.

\section{Methodische Abgrenzung der Eigenleistung}
Der Einsatz von generativen Sprachmodellen (Google Antigravity / LLM) beschränkte sich strikt auf Hilfs-, Beschleunigungs- und Syntax-Aufgaben:
\begin{itemize}
    \item Erstellung der zustandslosen Demo-Webanwendung (Mock-API und HTML-Dashboard), da die reine Anwendungsentwicklung nicht den wissenschaftlichen Kern der Arbeit darstellt.
    \item Bereitstellung von Python-Messschleifen zur Latenzerfassung.
    \item Syntaxvalidierung von Konfigurationsdateien (GitHub Actions, Kustomize, Terraform).
    \item Typografische Unterstützung bei der Formatierung von LaTeX-Tabellen.
\end{itemize}

Die Konzeption, die Formulierung der Forschungsleitfrage, die architektonischen Entwurfsentscheidungen, die Durchführung der Störungsexperimente auf AWS sowie die kritische Interpretation der Ergebnisse stellen originäre Eigenleistungen des Autors dar.

\section{Promptverzeichnis}
Nachfolgend sind die zehn zentralen Prompts aufgeführt, mit denen die unterstützenden Komponenten generiert und formatiert wurden:

\begin{enumerate}[leftmargin=*]
    \item \promptentry{FastAPI Demo-Backend}{„Schreibe mir ein minimales Python-FastAPI-Backend für einen Kubernetes-Testserver. Ich brauche Standard-Routen für \texttt{/health/live}, \texttt{/health/ready}, \texttt{/api/info} (liefert Hostname, IP und Uptime) sowie einen POST-Endpunkt \texttt{/lab/crash}, der den Prozess hart via \texttt{os.\allowbreak\_exit(1)} beendet.“}
    
    \item \promptentry{Web-Dashboard UI}{„Erstelle mir eine einfache, responsive HTML-Seite mit modernem Dark-Mode-CSS für die FastAPI-App. Die Seite soll per JavaScript-Fetch alle 1,5 Sekunden \texttt{/api/info} abfragen und den antwortenden Pod-Namen sowie die Latenz anzeigen, damit man Pod-Wechsel visuell nachvollziehen kann. Keine externen Frameworks wie React oder Tailwind.“}
    
    \item \promptentry{Gehärtetes Dockerfile}{„Erstelle mir ein schlankes Dockerfile auf Basis von \texttt{python:3.11-slim} für diese FastAPI-App. Achte darauf, dass die App unter einem unprivilegierten Benutzer (\texttt{non-root user}) läuft und ein Standard-HEALTHCHECK definiert ist.“}
    
    \item \promptentry{Lastgenerator-Skript}{„Schreibe mir ein Python-Skript mit \texttt{concurrent.\allowbreak futures.\allowbreak ThreadPoolExecutor}, das kontinuierlich 15 parallele HTTP-GET-Requests pro Sekunde gegen eine URL sendet, HTTP-Statuscodes zählt und am Ende die Perzentile $p_{50}$ und $p_{95}$ der Antwortzeiten als JSON ausgibt.“}
    
    \item \promptentry{GitOps-Drift Automatisierung}{„Erstelle mir ein Python-Skript, das über \texttt{kubectl} den Befehl \texttt{scale\allowbreak{} deployment\allowbreak{} -{}-replicas=0} absetzt und in einer Sekundenschleife den Status \texttt{status.\allowbreak sync.\allowbreak status} der Argo-CD-Application abfragt, um die exakte Zeitspanne bis zum Auftreten von \texttt{OutOfSync} sowie die Zeit bis zum Wiedererreichen von 2 Replikaten zu stoppen.“}
    
    \item \promptentry{GitHub Actions Workflow}{„Wie sieht die YAML-Syntax für eine GitHub-Actions-Pipeline aus, die bei einem Push auf \texttt{main} ein Docker-Image baut und sich mit dem automatischen \texttt{GITHUB\_TOKEN} an der GitHub Container Registry (\texttt{ghcr.io}) anmeldet?“}
    
    \item \promptentry{Kustomize-Overlays}{„Gib mir die Kustomize-Syntax für ein Base-Deployment (1 Replikat) und ein Overlay \texttt{overlays/aws}, das per Patch die Replica-Zahl auf 2 erhöht und eine Traefik-Ingress-Ressource hinzufügt.“}
    
    \item \promptentry{AWS Security Group Parameter}{„Ich erhalte bei Terraform den Fehler: \texttt{InvalidParameterValue:\allowbreak{} Value\allowbreak{} for\allowbreak{} parameter\allowbreak{} GroupDescription\allowbreak{} is\allowbreak{} invalid.\allowbreak{} Character\allowbreak{} sets\allowbreak{} beyond\allowbreak{} ASCII\allowbreak{} are\allowbreak{} not\allowbreak{} supported}. Woran liegt das in meiner \texttt{security-groups.tf}?“}
    
    \item \promptentry{k3s TLS-SAN-Konfiguration}{„Mein lokales \texttt{kubectl} wirft \texttt{x509:\allowbreak{} certificate\allowbreak{} is\allowbreak{} valid\allowbreak{} for\allowbreak{} 10.0.1.88,\allowbreak{} not\allowbreak{} 63.176.27.134}, wenn ich mich mit der Elastic IP verbinde. Wie konfiguriere ich \texttt{tls-san} in k3s nachträglich, damit das Zertifikat die öffentliche IP akzeptiert?“}
    
    \item \promptentry{LaTeX-Tabellengestaltung}{„Ich habe folgende Messergebnisse aus einem Störungsversuch: Zeitpunkte $t_0=0\text{ s}$ (Synced, 2 Pods), $t_2=2{,}6\text{ s}$ (OutOfSync), $t_4=11{,}3\text{ s}$ (Synced, 2 Pods). Formatiere mir diese Daten in sauberem LaTeX-Code unter Verwendung des \texttt{booktabs}-Pakets mit \texttt{\textbackslash toprule}, \texttt{\textbackslash midrule} und \texttt{\textbackslash bottomrule} passend für eine DIN-A4-Projektarbeit.“}
\end{enumerate}

    \newpage
    \chapter*{Eigenständigkeitserklärung}
\addcontentsline{toc}{section}{Eigenständigkeitserklärung}
Hiermit bestätige ich, dass ich die vorliegende Arbeit selbstständig verfasst und keine 
anderen Publikationen, Vorlagen und Hilfsmittel (z.B. künstliche Intelligenz) als die 
angegebenen benutzt habe. Alle Teile meiner Arbeit, die wortwörtlich oder dem Sinn nach 
anderen Werken entnommen sind, wurden unter Angabe der Quelle kenntlich gemacht. 
Gleiches gilt für von mir verwendete Internetquellen. Ich versichere, dass ich diese Arbeit 
oder nicht zitierte Teile daraus vorher nicht in einem anderen Prüfungsverfahren eingereicht 
habe. Mir ist bekannt, dass meine Arbeit zum Zwecke eines Plagiatsabgleichs mittels einer 
Plagiatserkennungssoftware auf eine ungekennzeichnete Übernahme von fremdem geistigen 
Eigentum sowie auf die Nutzung von künstlicher Intelligenz zur Texterstellung überprüft 
werden kann. Ich versichere, dass die elektronische Form meiner Arbeit mit der gedruckten 
Version identisch ist.


% \section*{Declaration of Originality}
% I hereby confirm that I have independently written this work and have not used any 
% publications, templates, or aids (e.g. artificial intelligence) other than those I have 
% indicated. All parts of my work which have been taken literally or correspondingly from other
% publications have been duly acknowledged. This also applies to Internet sources. I confirm
% that I have not previously submitted this work or any unquoted parts thereof in any other 
% examination procedure. I am aware that my work may be checked for plagiarism by means 
% of plagiarism recognition software, as well as for the use of artificial intelligence for text
% creation, in order to verify the integrity of its written content. I also confirm that the 
% electronic form is identical to the printed version. \\

\vspace{2cm}

\begin{center}
\begin{tabular}{p{0.45\textwidth}p{0.45\textwidth}}
    [Ort], \today & \\
    \hline
    % Place, Date & Signature \\
    Ort, Datum & Unterschrift \\
\end{tabular}
\end{center}

\end{onehalfspace}
\end{document}
